%%%%%%%%%%%%%%%%%%%%
% CHAPTER 4
% %%%%%%%%%%%%%%%%%%

% TODO: placeholder chapter image. Jack's leaning choice is a dramatic cone motif
% (thematically central to §4.8); replace before release. See docs/ch4_planning/brief.md.
\chapterimage{chapter_head_1.pdf} % Chapter heading image

\chapter{A Taste of Category Theory}

\section{What Are We Doing Here?}
We've spent three chapters unpacking three species ($\Z$, groups, and rings), and by now
we've probably noticed that the same tune keeps playing. Group homomorphisms and ring
homomorphisms look eerily similar; the direct product of two groups looks eerily like the
Cartesian product of two sets; the First Isomorphism Theorem for groups reads like the
same proof as its ring cousin, just with different nouns. Category theory is the language in
which that tune becomes a single melody.

This chapter is a \textit{taste}, not a first course in category theory. We want three
ideas in your hands (\textit{categories}, \textit{functors}, and \textit{natural
transformations}), plus one capstone, \textit{limits}, that ties every universal property we
meet into a single construction. Each idea will land on a MAWF example you already trust,
and if you want to go further,
\href{https://arxiv.org/abs/1803.05316}{Fong and Spivak's \textit{Seven Sketches}}
\cite{fong2018seven} or Roman's \textit{An Introduction to the Language of Category Theory}
\cite{roman2017introduction} are two natural next stops.

Category theory itself was invented in 1945 by
\href{https://en.wikipedia.org/wiki/Samuel_Eilenberg}{Samuel Eilenberg} and
\href{https://en.wikipedia.org/wiki/Saunders_Mac_Lane}{Saunders Mac Lane} in a paper titled
\textit{General Theory of Natural Equivalences}. Their actual problem was an annoyance in
algebraic topology, where the same map kept showing up in different disguises and they wanted
to say ``the same'' \textit{precisely}. To define a natural transformation they needed
functors; to define functors they needed categories; and so the language that now reorganizes
large parts of mathematics got built, a bit accidentally, from the top down. We'll walk the
same road, in the same order.

\noindent \textbf{Promises.} By the end of this chapter we will:
\begin{enumerate}[label=\roman*)]
    \item define categories and see that sets, groups, and rings live in them naturally;
    \item use the universal property of products to build the direct product of rings
    \textit{for free}, with no component-wise operation tables required;
    \item meet functors through forgetful functors and product-as-a-functor;
    \item meet natural transformations by comparing two different constructions on the same
    category;
    \item assemble all of the above into the single notion of a \textit{limit}, and watch
    products, equalizers, and even the kernels from Chapter 2 fall out as special cases.
\end{enumerate}

\noindent A friendly warning: Chapter 4 is genuinely more abstract than the chapters before
it. We will manage the abstraction out loud. Whenever we climb the ladder, we'll come back
down and cash out what the climb bought us in terms of things you should by now be quite
familiar with.

Speaking of which: \hyperref[note:ch1-gcd-teaser]{back in Chapter 1 we posed a question to
the CT-aware reader}. Does the definition of $GCD(a,b)$ look familiar? Can you think of a
category in which GCD \textit{is} a categorical product? We are now a few sections away from
being able to answer that, and we will in \S4.4. First, we need categories.

\section{Categories}
% TODO: see docs/ch4_planning/ch4_outline_v3.md §4.2 (Chunks 4.2a--4.2c)

\section{Set as Home Base: Terminal and Initial Objects}
% TODO: see docs/ch4_planning/ch4_outline_v3.md §4.3 (Chunks 4.3a--4.3b)

\section{Products via Universal Property}
% TODO: see docs/ch4_planning/ch4_outline_v3.md §4.4 (Chunks 4.4a--4.4d)

\section{Coproducts and the Duality We've Been Using}
% TODO: see docs/ch4_planning/ch4_outline_v3.md §4.5 (Chunks 4.5a--4.5b)

\section{Functors}
% TODO: see docs/ch4_planning/ch4_outline_v3.md §4.6 (Chunks 4.6a--4.6c)

\section{Natural Transformations}
% TODO: see docs/ch4_planning/ch4_outline_v3.md §4.7 (Chunks 4.7a--4.7c)

\section{Limits}
% TODO: see docs/ch4_planning/ch4_outline_v3.md §4.8 (Chunks 4.8a--4.8c)

\section{Where to Go From Here}
% TODO: see docs/ch4_planning/ch4_outline_v3.md §4.9 (Chunk 4.9)
